This thesis examines whether a practical Linux \gls{cli} assistant can be built
with small local \glspl{llm} under offline and hardware constraints.
The thesis proposes shAI, a modular \gls{rag}-based prototype composed of a
command-vs.-query classifier, a document retriever that indexes Linux man pages,
and a summarizer that generates concise answers from retrieved passages.

The classifier was tested on the balanced\\
\texttt{classifier-100} dataset; accuracy ranged from 50\% to 88\%, with the
best model requiring 2.9s per query.
The retriever accuracy was evaluated on a collections of 10 and 100 documents 
(\texttt{sample-10} and \texttt{sample-100} respectively)
with different retrieval approaches. Document-level retrieval approach has reached
top-5 document hits from 86.3\% to 92.0\% with MRR ranging from 0.724 to 0.806
on the 100-document collection, generating huge output sizes (up to 120K
tokens). To address that, document-level with passage-level reranking and
passage-level retrieval approaches were implemented, showing much more compact
genrated output sizes (~5K tokens) and achieving precision.

The application latency was evaluated on \texttt{sample-10},
\texttt{sample-100}, and\\ \texttt{sample-1000}, showing that passage-level
retrieval and proposed document-level with passage-level reranking approaches
achieved the best overall performance, based on the size of chosen embedding
model, with latencies under 2s for search and under 10s for RAG.

The results demonstrate that retrieval-first architectures are practical for
local \gls{cli} assistance, whereas agentic open-ended approaches remain
unreliable on small models, with the best choice for the retrieval technique
being the passage-level retrieval. The thesis provides an open-source prototype and
empirical benchmarks that guide future work on constrained task design.
